problem:  
Let $\pi:(E,g)\to (M,h)$ be a Riemannian submersion whose total space $E$ is a smooth sphere bundle with fiber $S^k$ and compact base $M$.  Recall that the O’Neill curvature equations decompose the sectional curvatures of $E$ into horizontal, vertical, and mixed parts.  

Assume that:  
1. all fibers are totally geodesic,  
2. the horizontal distribution is *fat* in the sense of Weinstein—that is, for every nonzero horizontal vector $X$ and nonzero vertical vector $V$, the **vertizontal curvature** satisfies $K(X,V)>0$.  

Question:  
Under these hypotheses, is it true that there exists a smooth deformation of the metric $g$, keeping the submersion structure (fibers totally geodesic and horizontally/vertically compatible), such that the resulting metric on $E$ has strictly positive sectional curvature on *all* 2‑planes?  

Equivalently, do fat sphere bundles always admit metrics of strictly positive sectional curvature compatible with the given submersion structure?  

Why is it a "good" problem:  
This formulation isolates a precise and long‑standing gap in the theory of Riemannian submersions: while Weinstein’s fatness condition guarantees positivity of vertizontal curvatures and rules out obvious topological obstructions, it has remained unclear whether fatness suffices—or can be deformed to suffice—for full sectional positivity on the total space.  The problem bridges classical bundle topology and modern global Riemannian geometry, connecting O’Neill’s submersion formulas, canonical metric variations, and the topological classification of sphere bundles.  

Its statement is simple but its implications are vast: a positive resolution would present a systematic method to construct new compact manifolds of positive sectional curvature directly from fat bundle data; a negative resolution would reveal hidden obstructions in the curvature coupling between horizontal and vertical directions.  Thus this question sits at the intersection of geometric analysis, differential topology, and curvature theory, embodying the characteristics of a “good” research problem—conceptually clear, deeply interconnected, and of central geometric importance.